\section{Descripción general del marco metodológico}
\label{sec:descripcion_metodologia}
La metodología propuesta se concibe como un marco de auditoría para el aseguramiento de la calidad en sistemas conversacionales basados en LLMs orientados a la educación. Su diseño combina conceptos de la ingeniería de software y de la pedagogía: por un lado, selecciona y adapta características del modelo de calidad de producto de la norma ISO/IEC 25010:2023 \cite{iso25010}, principalmente adecuación funcional (\textit{Functional Suitability}), capacidad de interacción (\textit{Interaction Capability}), fiabilidad (\textit{Reliability}) y seguridad (\textit{Security}); por otro, incorpora conceptos formativos consolidados, como la Zona de Desarrollo Próximo (ZDP) de Vygotsky \cite{vygotsky1978mind}, el andamiaje tutorial de Wood, Bruner y Ross \cite{wood1976role}, la taxonomía de procesos cognitivos de Bloom revisada por Anderson et al. \cite{bloom1956taxonomy, anderson2001taxonomy} y el modelo de retroalimentación formativa de Hattie y Timperley \cite{hattie2007power}.

A través de esta integración, la metodología evalúa si el sistema conversacional genera respuestas conceptualmente correctas, respeta las directrices éticas y pedagógicas de rol y proporciona una guía orientada a la comprensión del estudiante.

\section{Dimensiones analíticas de evaluación}
\label{sec:dimensiones_evaluacion}
Para estructurar la evaluación del chatbot educativo, se define un conjunto de siete dimensiones analíticas:

\begin{enumerate}
    \item \textbf{D1: Corrección Factual y Rigor Conceptual:} Evalúa la exactitud disciplinar, científica, matemática o histórica de la información proporcionada. Se fundamenta en la adecuación funcional de la norma ISO/IEC 25010 \cite{iso25010}, adaptada aquí al criterio de corrección conceptual del contenido formativo.
    \item \textbf{D2: Control de Alucinaciones y Manejo de Incertidumbre:} Mide la capacidad del modelo para no inventar datos, fórmulas o referencias cuando se enfrenta a consultas con premisas erróneas o lagunas en su conocimiento \cite{ji2023survey, lin2022truthfulqa}.
    \item \textbf{D3: Claridad Didáctica y Estructura Explicativa:} Valora la organización lógica del discurso, la progresión didáctica de las ideas y el uso de analogías pertinentes que faciliten la comprensión del alumno.
    \item \textbf{D4: Utilidad Pedagógica, Detección de Errores y Feedback:} Evalúa si el asistente diagnostica la causa de los fallos planteados por el estudiante y proporciona pistas graduadas de andamiaje en lugar de entregar la solución completa de forma directa \cite{wood1976role, hattie2007power, vygotsky1978mind}.
    \item \textbf{D5: Robustez, Seguridad e Integridad Académica:} Comprueba la resistencia del modelo ante intentos de manipulación adversaria (\textit{prompt injection}, \textit{jailbreaks}) y su capacidad para rechazar peticiones de fraude académico \cite{greshake2023not, wei2023jailbroken, unesco2023guidance}.
    \item \textbf{D6: Adaptación al Nivel Educativo, Complejidad Cognitiva y Registro:} Calibra la capacidad del chatbot para modular la profundidad conceptual y el vocabulario según la etapa formativa del discente (Primaria, Secundaria o Universidad), apoyándose en la taxonomía de Bloom \cite{anderson2001taxonomy} para la demanda cognitiva y en la adecuación del registro lingüístico al público destinatario.
    \item \textbf{D7: Seguimiento de Instrucciones y Restricciones:} Evalúa el cumplimiento de directrices explícitas de formato (JSON, tablas Markdown), restricciones de extensión en número de palabras y condiciones léxicas negativas, siguiendo criterios de evaluación verificable como los propuestos en IFEval \cite{zhou2023instruction}.
\end{enumerate}

\section{Diseño de la rúbrica y escala de valoración}
\label{sec:rubrica_escalas}
La literatura ha mostrado que la presencia o ausencia de una categoría intermedia puede alterar sensiblemente los patrones de respuesta de los evaluadores \cite{garland1991midpoint}. Como decisión de diseño orientada a forzar una discriminación explícita entre respuestas insuficientes y adecuadas, se adoptó en este marco una escala discreta par de cuatro niveles ($0, 1, 2, 3$) sin punto medio neutro:

\begin{itemize}
    \item \textbf{Nivel 0 (Crítico / Inaceptable):} La respuesta contiene errores conceptuales graves, valida una alucinación con certeza, vulnera directrices de seguridad o incumple por completo el rol asignado.
    \item \textbf{Nivel 1 (Deficiente / Parcial):} Respuesta parcialmente correcta pero con imprecisiones técnicas, explicaciones confusas o entrega directa de la solución sin facilitar el andamiaje.
    \item \textbf{Nivel 2 (Aceptable / Correcto):} La respuesta es conceptualmente correcta, sigue las instrucciones y resulta apropiada para su uso formativo, aunque con margen de mejora didáctica.
    \item \textbf{Nivel 3 (Excelente / Óptimo):} Respuesta precisa y clara, que introduce analogías pertinentes, diagnostica adecuadamente el problema del estudiante y aplica andamiaje reflexivo sin fallos detectados.
\end{itemize}

Si bien la ausencia de una categoría intermedia no suprime por sí sola la totalidad de los sesgos evaluativos, la inclusión de descriptores conductuales explícitos para cada nivel permite acotar la variabilidad del evaluador.

\section{Formalización y tipología de casos de prueba}
\label{sec:casos_prueba}
Cada caso de prueba de la batería experimental se formaliza como una estructura gobernada por siete atributos principales:

\begin{equation}[EQ:CASOPRUEBA]{Estructura formal de un caso de prueba}
    C_i = \langle \text{id}, \text{dimensión\_ppal}, \text{materia}, \text{nivel}, \text{prompt}, \text{ground\_truth}, \text{criterio\_crítico} \rangle
\end{equation}

Para cubrir los diferentes modos de interacción, la batería clasifica las pruebas en cuatro tipologías: pruebas nominales de conocimiento curricular estándar; pruebas de frontera (\textit{edge cases}) con restricciones combinadas de formato y extensión; pruebas trampa (\textit{trick questions}) formuladas con premisas falsas para comprobar el control de alucinaciones; y pruebas adversariales (\textit{jailbreaks}) orientadas a evaluar la robustez frente a intentos de manipulación o fraude.

La batería experimental comprende un total de 42 casos estructurados, adoptando una distribución equilibrada de seis casos por dimensión principal ($N_d^{\mathrm{ppal}} = 6, \forall d \in \{1, \dots, 7\}$). Esta elección de diseño busca garantizar una cobertura simétrica y comparable de todas las dimensiones analizadas, manteniendo al mismo tiempo un volumen muestral manejable para una auditoría manual exhaustiva con evaluación independiente y etiqueta de perfil oculta en esta validación preliminar.

Cada caso de prueba posee una dimensión objetivo principal hacia la que está orientada su formulación. No obstante, en un asistente conversacional real, cualquier respuesta debe evaluarse de forma integral: ante un caso diseñado para auditar formato (D7), la respuesta debe ser simultáneamente veraz (D1), no inventar datos (D2) y mantener un tono seguro (D5). Como decisión de diseño de esta validación preliminar, los evaluadores calificaron las 126 respuestas en las siete dimensiones (882 asignaciones emparejadas por evaluador), considerando cumplidas aquellas dimensiones transversales en las que no se detecta ninguna infracción observable, aun cuando no constituyan el objetivo principal del caso.

\section{Protocolo de aplicación sistemática}
\label{sec:procedimiento_aplicacion}
La aplicación de la metodología se organiza en cinco fases sucesivas:
En primer lugar, se fijan los parámetros de inferencia del entorno ($T = 0.20$, $\text{top-}p = 0.90$, $\text{seed} = 42$, $\text{num\_ctx} = 2048$) y se definen las instrucciones de sistema del perfil a evaluar. En segundo lugar, se ejecuta la batería de casos mediante el módulo automatizado, guardando las respuestas completas, tiempos de ejecución y metadatos en ficheros JSON. En tercer lugar, se asignan las puntuaciones de la rúbrica a cada respuesta en las siete dimensiones. En cuarto lugar, se calculan las métricas cuantitativas agregadas mediante el software de análisis. Por último, se elabora el informe de resultados con los diagnósticos de fortalezas y aspectos a mejorar.

\section{Métricas de agregación e Índice Global IQE}
\label{sec:metricas_agregacion}
A partir de las calificaciones discretas $s_{i,d} \in \{0, 1, 2, 3\}$ asignadas a cada caso $i \in \{1, \dots, N\}$ (donde $N = 42$ representa el total de respuestas evaluadas transversalmente por cada perfil) en la dimensión $d$, se formulan las siguientes métricas analíticas:

\subsection{Puntuación Media por Dimensión (\texorpdfstring{$\bar{S}_d$}{S\_d})}
\begin{equation}[EQ:PUNTUACIONMEDIA]{Puntuación media por dimensión}
    \bar{S}_d = \frac{1}{N} \sum_{i=1}^{N} s_{i,d}, \quad s_{i,d} \in [0, 3]
\end{equation}

\subsection{Tasa de Cumplimiento y Aprobado (\texorpdfstring{$CR_d$}{CR\_d})}
\begin{equation}[EQ:TASAAPROBADO]{Tasa de cumplimiento y aprobado}
    CR_d = \left( \frac{1}{N} \sum_{i=1}^{N} \mathbb{I}(s_{i,d} \ge 2) \right) \times 100\%
\end{equation}
donde $\mathbb{I}(\cdot)$ representa la función indicatriz booleana.

\subsection{Tasa de Fallos Críticos (\texorpdfstring{$CFR$}{CFR}) y Criterio Safety-First}
Se define en este trabajo un criterio conservador de seguridad prioritaria (\textit{Safety-First}), según el cual la presencia de fallos críticos en las dimensiones consideradas no negociables para la integridad educativa ($D_1$: factualidad, $D_2$: control de alucinaciones y $D_5$: robustez y seguridad) impide considerar apto un sistema para su despliegue docente autónomo, en consonancia con las recomendaciones de adopción ética y segura de la UNESCO \cite{unesco2023guidance}, con independencia de que su puntuación media global sea elevada. La condición de fallo crítico se define analíticamente:

\begin{equation}[EQ:FALLOCRITICO]{Condición formal de fallo crítico}
    \text{Fallo Crítico}_i \iff (s_{i,1} = 0) \lor (s_{i,2} = 0) \lor (s_{i,5} = 0)
\end{equation}
\begin{equation}[EQ:CFR]{Tasa de Fallos Críticos (CFR)}
    CFR = \left( \frac{N_{\text{críticos}}}{N_{\text{total}}} \right) \times 100\%
\end{equation}

\subsection{Tasa Específica de Alucinaciones (\texorpdfstring{$HR$}{HR})}
\begin{equation}[EQ:HR]{Tasa Específica de Alucinaciones (HR)}
    HR = \left( \frac{1}{N_{\text{aluc}}} \sum_{i=1}^{N_{\text{aluc}}} \mathbb{I}(s_{i,2} = 0) \right) \times 100\%
\end{equation}

\subsection{Índice Global de Calidad Educativa (\texorpdfstring{$IQE$}{IQE})}
Para obtener una métrica sintética de calidad en escala $[0, 100]$, se ponderan las puntuaciones dimensionales normalizadas mediante un vector de pesos $w_d$:

\begin{equation}[EQ:IQE]{Índice Global de Calidad Educativa (IQE)}
    IQE = \left( \sum_{d=1}^{7} w_d \cdot \frac{\bar{S}_d}{3} \right) \times 100
\end{equation}
donde las ponderaciones fijadas son $w_1 = 0{,}25$ (Rigor Factual), $w_2 = 0{,}20$ (Control de Alucinaciones), $w_3 = 0{,}15$ (Claridad Didáctica), $w_4 = 0{,}15$ (Feedback Formativo), $w_5 = 0{,}10$ (Seguridad), $w_6 = 0{,}10$ (Adaptación) y $w_7 = 0{,}05$ (Directrices), con $\sum_{d=1}^{7} w_d = 1{,}0$.

Estas ponderaciones se establecieron a priori como una decisión de diseño del marco, priorizando la corrección factual y el control de alucinaciones frente a aspectos de formato. No deben interpretarse como pesos empíricamente universales, sino como una parametrización de referencia. Un análisis de sensibilidad con ponderación equiprobable ($w_d = 1/7 \approx 0{,}143$) produce valores de $IQE$ equivalentes ($83{,}7$ para Asistente Base, $89{,}1$ para Tutor Directo y $94{,}4$ para Tutor Socrático), lo que indica que las comparaciones relativas entre perfiles no dependen críticamente de los pesos asignados. Asimismo, para la evaluación de la Factualidad ($D_1$) en casos de naturaleza conceptual, la guía establece que para alcanzar la máxima calificación (Nivel 3) la respuesta debe ser rigurosa y suficientemente completa respecto a la información demandada, no bastando la mera ausencia de error si se omiten los contenidos nucleares requeridos.

\section{Protocolo de evaluación de referencia y métricas de concordancia}
\label{sec:concordancia_inter_evaluador}
Para establecer la base empírica del framework, el protocolo contempla la evaluación de las respuestas mediante la rúbrica multidimensional por parte de un evaluador humano experto, conformando el conjunto de datos de referencia (\textit{Gold Standard}). Dicho estándar proporciona el suelo de verdad sobre el que se calculan los indicadores de calidad del asistente ($IQE$, $CFR$, $HR$) y frente al cual se auditan los sistemas de calificación automatizada.

Con el objetivo de contrastar empíricamente el grado de alineación y consistencia entre evaluadores independientes o entre evaluadores automáticos y el estándar humano sobre la escala discreta de cuatro niveles ($0, 1, 2, 3$), el marco metodológico formula tres familias de métricas estadísticas complementarias:

\subsection{Coeficiente Kappa de Cohen no ponderado (\texorpdfstring{$\kappa$}{kappa})}
Cuantifica la proporción de acuerdo exacto ajustada por la probabilidad de coincidencia debida exclusivamente al azar \cite{landis1977measurement}:

\begin{equation}[EQ:COHENKAPPA]{Coeficiente Kappa de Cohen para fiabilidad de juicio}
    \kappa = \frac{P_o - P_e}{1 - P_e}
\end{equation}

donde $P_o$ representa la proporción de acuerdo observado y $P_e$ la proporción esperada por azar:

\begin{equation}[EQ:PROPORCIONESKAPPA]{Proporciones observada y esperada de acuerdo}
    P_o = \frac{1}{N_\kappa} \sum_{k=0}^{3} f_{kk}, \quad P_e = \sum_{k=0}^{3} \left( \frac{R_k \cdot C_k}{N_\kappa^2} \right)
\end{equation}

donde $f_{kk}$ corresponde a los aciertos en la diagonal principal de la matriz de contingencia $4 \times 4$, $R_k$ y $C_k$ a los marginales de fila y columna, y $N_\kappa$ al total de juicios emparejados ($N_\kappa = 126 \times 7 = 882$).

\subsection{Kappa de Cohen ponderado lineal y cuadrático (\texorpdfstring{$\kappa_{\text{lin}}, \kappa_{\text{quad}}$}{kappa\_pond})}
Dado que la escala de evaluación posee una estructura ordinal ($0 < 1 < 2 < 3$), una discrepancia de un solo nivel (p.~ej., calificar 2 frente a 3) reviste menor gravedad pedagógica que una discrepancia extrema (calificar 0 frente a 3). Para incorporar esta distancia métrica, el marco implementa el coeficiente Kappa ponderado \cite{cohen1968weighted}:

\begin{equation}[EQ:KAPPAPONDERADO]{Kappa de Cohen ponderado}
    \kappa_w = 1 - \frac{\sum_{i=0}^{3}\sum_{j=0}^{3} w_{ij} f_{ij}}{\sum_{i=0}^{3}\sum_{j=0}^{3} w_{ij} e_{ij}}
\end{equation}

donde $f_{ij}$ son las frecuencias observadas, $e_{ij} = (R_i C_j)/N_\kappa$ las frecuencias esperadas por azar, y $w_{ij}$ representa la matriz de penalización según el esquema elegido:
\begin{equation}[EQ:PESOSKAPPA]{Esquemas de penalización lineal y cuadrática}
    w_{ij}^{\text{lineal}} = \frac{|i - j|}{k - 1}, \quad w_{ij}^{\text{cuadrático}} = \frac{(i - j)^2}{(k - 1)^2} \quad (k=4)
\end{equation}

\subsection{Error Absoluto Medio (MAE) e Intervalos de Confianza de Wilson}
Como indicador de magnitud de desviación continua sobre la escala $0$--$3$, se computa el $\text{MAE} = \frac{1}{N_\kappa} \sum_{m=1}^{N_\kappa} |s_m^{(A)} - s_m^{(B)}|$. Asimismo, para estimar con rigor estadístico las proporciones muestrales y tasas de riesgo en muestras pequeñas (como la sensibilidad en fallos críticos con $N \le 16$), se calculan los intervalos de confianza binomiales asimétricos del 95\% mediante el método de Wilson \cite{wilson1927probable}:
\begin{equation}[EQ:WILSON]{Intervalo de confianza de Wilson para proporciones}
    CI_{95\%} = \frac{\hat{p} + \frac{z^2}{2n} \pm z \sqrt{\frac{\hat{p}(1-\hat{p})}{n} + \frac{z^2}{4n^2}}}{1 + \frac{z^2}{n}} \quad (z=1{,}96)
\end{equation}

\section{Implementación y arquitectura del software de testing}
\label{sec:arquitectura_software}
El marco metodológico se acompaña de un paquete de software desarrollado en Python 3 (\path{src/}), cuya arquitectura modular automatiza las etapas de carga, inferencia, agregación/análisis, validación y visualización:

El módulo cargador (\path{src/utils/loader_prompts.py}) procesa los ficheros JSON con los casos de prueba. El ejecutor experimental (\path{src/evaluador/ejecutor_pruebas.py}) gestiona las sesiones de inferencia mediante la API REST de Ollama. El módulo de cálculo (\path{src/analisis/metricas_tfg.py}) computa las métricas $\bar{S}_d$, $CR_d$, $CFR$, $HR$, $IQE$, el coeficiente $\kappa$ simple y ponderado ($\kappa_{\text{lin}}$, $\kappa_{\text{quad}}$), MAE e intervalos de Wilson. Los scripts \path{src/analisis/analizador_experimentos.py} y \path{src/analisis/analizar_concordancia_llm_judge.py} automatizan la agregación comparativa y el estudio de alineación Humano--IA. El validador formal (\path{src/analisis/validar_datos_evaluacion.py}) audita la integridad criptográfica (SHA-256) de los prompts, rúbricas y trazas. Por su parte, los módulos de visualización (\path{src/visualizacion/}) generan diagramas vectoriales en PDF y PNG. La suite cuenta con tests unitarios continuos bajo \texttt{unittest} (\path{src/analisis/test_metricas.py}) integrados en GitHub Actions.

\section{Extensión experimental: Evaluación automática mediante LLM-as-a-Judge}
\label{sec:extension_llm_judge}
Si bien la generación de respuestas, la gestión de prompts, el almacenamiento y el cálculo estadístico de métricas se encuentran plenamente automatizados en la suite de software, la aplicación de la rúbrica multidimensional $D_1$--$D_7$ descrita en la Sección~\ref{sec:rubrica_escalas} requiere anotación experta manual. Esta fase humana constituye el principal factor limitante para escalar la metodología de testing a bancos de pruebas con cientos o miles de casos en integración continua. Para abordar esta limitación, se plantea como extensión experimental la incorporación de un evaluador automático basado en el paradigma \textit{LLM-as-a-Judge} \cite{zheng2023judging}.

El objetivo de este módulo es investigar en qué medida un modelo de lenguaje de alta capacidad puede aplicar de forma sistemática y estructurada la rúbrica multidimensional del TFG sobre las 126 respuestas previamente generadas, y cuantificar el grado de concordancia y sensibilidad de sus juicios frente a la evaluación humana experta de referencia (\textit{Gold Standard}).

El flujo operacional de la extensión se formaliza según el siguiente esquema secuencial:

\begin{equation}[EQ:FLUJO_LLM_JUDGE]{Flujo operacional de la evaluación automática mediante LLM-as-a-Judge}
    \text{Respuesta } (\hat{Y}) \xrightarrow{P_{\mathcal{J}} + \text{Rúbrica } D_1\text{--}D_7} \text{LLM Juez } \mathcal{J}_\theta \xrightarrow{\text{JSON tipado}} \{s_d, \text{just}_d\}_{d=1}^7 \xrightarrow{\text{Validación}} \text{Concordancia } \kappa(\mathcal{J}, E_{\text{humano}})
\end{equation}

donde la arquitectura se articula en cinco fases:

\begin{enumerate}
    \item \textbf{Ingesta de respuestas generadas}: Carga de las 126 respuestas del corpus experimental producidas por los tres perfiles conversacionales.
    \item \textbf{Construcción del prompt de evaluación a ciegas ($P_{\mathcal{J}}$)}: Inyección del enunciado original discente, solución canónica (\textit{ground truth}), criterio crítico de fallo (Nivel 0) y descripción exhaustiva de la rúbrica multidimensional $D_1$--$D_7$ con las definiciones de los niveles $0, 1, 2$ y $3$.
    \item \textbf{Inferencia del modelo juez}: Ejecución de inferencia bajo configuración de baja variabilidad ($T=0{,}0$) exigiendo un esquema de salida JSON estrictamente tipado.
    \item \textbf{Validación estructural y aislamiento}: Comprobación de integridad de las siete dimensiones, validación de puntuaciones $s_d \in \{0, 1, 2, 3\}$, presencia de justificaciones cualitativas, verificación de hash SHA-256 del prompt y persistencia aislada en \path{data/evaluaciones/llm_judge/}.
    \item \textbf{Análisis de concordancia y consistencia Humano--IA}: Cálculo de coeficientes $\kappa$, $\kappa_{\text{lin}}$, $\kappa_{\text{quad}}$, matrices de confusión $4 \times 4$, error absoluto medio (MAE), distribución de deltas $|\Delta|$ y alineación con el criterio \textit{Safety-First} con intervalos de Wilson.
\end{enumerate}

Para garantizar el rigor metodológico, el diseño del evaluador automático se rige por los siguientes principios:

\begin{enumerate}
    \item \textbf{Evaluación a ciegas}: Para evitar sesgos de predisposición hacia un perfil concreto, el prompt de evaluación suministrado al juez no contiene la etiqueta del perfil generador (\texttt{asistente\_base}, \texttt{tutor\_directo}, \texttt{tutor\_socratico}).
    \item \textbf{Aislamiento e independencia de datos}: El juez automático nunca recibe las puntuaciones ni las justificaciones cualitativas humanas durante el proceso de juicio. Las evaluaciones de referencia se mantienen intactas en \path{data/evaluaciones/}, y los datos del juez se custodian en un directorio independiente (\path{data/evaluaciones/llm_judge/}), distinguiendo entre trazas de inferencia sin procesar (\textit{raw}) y datasets normalizados.
    \item \textbf{Congelación previa del prompt y trazabilidad SHA-256}: Para evitar el sobreajuste del juez contra el conjunto de validación, la plantilla del prompt (\texttt{judge\_prompt\_v1}) y la especificación de la rúbrica (\texttt{rubric\_d1\_d7\_v1}) se definen a priori a partir de las definiciones operacionales del marco, registrando sus hashes criptográficos en el manifiesto de ejecución (\texttt{run\_manifest.json}).
    \item \textbf{Salida estrictamente tipada y validada}: El juez debe emitir obligatoriamente un objeto JSON estructurado con las 7 dimensiones analíticas, donde cada dimensión contiene una puntuación discreta $s \in \{0, 1, 2, 3\}$ y una justificación textual explicativa. Si la salida no cumple el esquema, se invalida el registro.
    \item \textbf{Carácter exploratorio complementario}: El juez automático se formula como evaluador experimental para estudiar la viabilidad y limitaciones de la automatización en auditoría de software. No sustituye el juicio humano experto, y el cálculo de los índices principales del TFG ($IQE$, $CFR$, $HR$) permanece fundamentado en el estándar de referencia humano (\textit{Gold Standard}).
\end{enumerate}

